\documentclass[12pt]{article}
\usepackage{amsmath, amsthm, amssymb}
\usepackage{geometry}
\geometry{margin=1.4in}
\usepackage{hyperref}
\usepackage{setspace}
\setstretch{1.4}

\newtheorem{theorem}{Theorem}
\theoremstyle{definition}
\newtheorem{definition}[theorem]{Definition}
\theoremstyle{remark}
\newtheorem{remark}[theorem]{Remark}

\title{The Universal Template:\\
How to Apply Displacement Framework to Any Domain}
\author{Diego Rincón\\
\small Independent}
\date{}

\begin{document}

\maketitle

\begin{abstract}
The displacement framework is not domain-specific. It applies everywhere: medicine, economics, politics, psychology, ecology, software, art, relationships, education, energy, physics. This paper gives the template. Given any domain, you can derive the displacement analysis. The variance is unlimited. The pattern is singular.
\end{abstract}

\section{The Universal Structure}

Every system has the same anatomy:
\begin{enumerate}
\item \textbf{Intended ground state} $S_0$: what the system was designed for, what we want
\item \textbf{Displacement} $\xi$: the pressure that pushes the system away
\item \textbf{Cost of displacement} $D(\xi)$: how much energy/resources the deviation requires to maintain
\item \textbf{Attractor} $S^*$: the stable state the system actually occupies
\item \textbf{Return cost} $C_{\text{return}}$: how much it costs to go back to $S_0$
\end{enumerate}

This structure is present in every domain.

\section{The Template for Any Domain}

\begin{definition}[Domain instantiation of displacement framework]
Given a domain (medicine, economics, politics, art, ecology, software, etc.):

\begin{enumerate}
\item Name the intended ground state $S_0$. What was this system designed for? What should it do?

\item Identify the pressure $\xi$ that pushes away. What force causes displacement? What makes deviating cheaper than staying?

\item Measure the cost $D(\xi)$. How much energy does the system spend maintaining the displacement?

\item Locate the actual attractor $S^*$. Where does the system actually settle? What is it optimized for instead?

\item Calculate the return barrier $C_{\text{return}}$. How much would it cost to go back to $S_0$?

\item Diagnose irreversibility. Has the system built infrastructure around $S^*$ that makes return impossible (or prohibitively expensive)?

\item Propose intervention. What would force a return? What would change the cost landscape?
\end{enumerate}
\end{definition}

\section{Examples Across Domains}

\subsection{Medicine}

\textbf{Ground state} $S_0$: Health. The body functioning as designed.

\textbf{Pressure} $\xi$: Pathogen, toxin, stress, genetics.

\textbf{Cost} $D(\xi)$: Energy spent fighting the pressure (fever, inflammation, immune response).

\textbf{Attractor} $S^*$: Chronic disease state. The body settled into a new stable configuration: inflamed, dysregulated, but stable.

\textbf{Return barrier}: High. The body has built tolerance to the new state. Inflammation has become part of the immune baseline. Returning to $S_0$ means dismantling this infrastructure.

\subsection{Economics}

\textbf{Ground state} $S_0$: Flourishing. Sustainable prosperity, distributed wealth, full employment.

\textbf{Pressure} $\xi$: Market incentives. Profit maximization.

\textbf{Cost} $D(\xi)$: Energy spent extracting value, externalizing costs.

\textbf{Attractor} $S^*$: Inequality. Concentrated wealth is the cheap stable state.

\textbf{Return barrier}: Prohibitive. Requires removing subsidies, pricing carbon, international coordination.

\subsection{Relationships}

\textbf{Ground state} $S_0$: Secure attachment. Two people meeting, understanding, being known.

\textbf{Pressure} $\xi$: Conflict, misunderstanding, competing needs.

\textbf{Cost} $D(\xi)$: Energy spent managing conflict, hiding, protecting.

\textbf{Attractor} $S^*$: Distant stability. Two people who have learned to coexist without intimacy.

\textbf{Return barrier}: High. Trust is broken. Returning to $S_0$ requires vulnerability from both people simultaneously—expensive.

\section{The Variance is Unlimited}

Every domain follows this pattern. But the specific instantiations are infinite:

\begin{itemize}
\item Sports: athletic potential vs. injury-avoidant sedentary comfort
\item Cooking: flavor and nutrition vs. fast, cheap, reproducible
\item Architecture: beauty and community vs. efficient profit-maximal building
\item Gardening: thriving ecosystem vs. monoculture lawn
\item Spirituality: meaning and connection vs. comfortable nihilism
\item Language: precision and expressiveness vs. cliché and shortcuts
\item Scientific research: discovery vs. publishable findings
\item Parenting: autonomy-building vs. control-maximizing
\end{itemize}

The list is infinite because every system has structure, pressure, and multiple stable states.

\section{The Insight}

Most interventions fail because they treat the symptom, not the displacement. A person with chronic inflammation takes anti-inflammatories instead of addressing why the body is stuck in inflammation. A country tries to fix inequality with redistribution instead of removing the incentive that creates concentration.

Successful intervention requires:
\begin{enumerate}
\item Acknowledging that $S^*$ is stable and the system is working correctly
\item Recognizing that return is expensive
\item Changing the cost structure so that $S_0$ becomes cheaper than $S^*$
\item Being willing to dismantle the infrastructure that supports $S^*$
\end{enumerate}

Without these, the system returns to its attractor within weeks.

\section{For the Reader}

Pick any domain you know. Apply the template:
\begin{enumerate}
\item What was intended? ($S_0$)
\item What pushes away? ($\xi$)
\item Where is it stuck? ($S^*$)
\item Why can't it return? ($C_{\text{return}}$)
\item What would make return cheaper? (intervention)
\end{enumerate}

You have now done displacement analysis. The pattern is universal. The applications are infinite.

The framework is not a collection of papers. It is a way of seeing.

\end{document}
